% ==============================================================================
% preamble/typography.tex — การตั้งค่าฟอนต์ การเยื้อง และระยะบรรทัด
% ตามมาตรฐาน thai-academic-latex-formatting
% ==============================================================================

% 1. ฟอนต์ภาษาไทยและอังกฤษ Sarabun (บรรจุในโฟลเดอร์ fonts/ แบบพกพาสมบูรณ์)
\XeTeXlinebreaklocale "th"
\XeTeXlinebreakskip = 0pt plus 1pt

\defaultfontfeatures{Mapping=tex-text}
\setmainfont{THSarabun.ttf}[
    Path = fonts/,
    BoldFont = THSarabun-Bold.ttf,
    ItalicFont = THSarabun-Italic.ttf,
    BoldItalicFont = THSarabun-BoldItalic.ttf,
    Script = Thai,
    Scale = 1.15
]

\newfontfamily\thaifont{THSarabun.ttf}[
    Path = fonts/,
    BoldFont = THSarabun-Bold.ttf,
    ItalicFont = THSarabun-Italic.ttf,
    BoldItalicFont = THSarabun-BoldItalic.ttf,
    Script = Thai,
    Scale = 1.15
]

\setmonofont{Courier}[Scale=0.92]

% 2. ระยะห่างระหว่างบรรทัด (ประมาณ 1.18 เท่า ตามข้อกำหนด)
\setstretch{1.18}

% 3. การเยื้องย่อหน้า (Paragraph Indentation)
% ข้อกำหนดเฉพาะ: เยื้อง 3em ไม่มีระยะห่างพิเศษระหว่างย่อหน้า (parskip = 0pt)
\setlength{\parindent}{3em}
\setlength{\parskip}{0pt}

% 4. รายการ (Lists Formatting)
% ไม่ให้การเยื้อง 3em ของย่อหน้าส่งผลกระทบต่อรายการซ้ำซ้อน (Retain distinct list indentation)
\setlist[itemize]{
    leftmargin=2.8em,
    itemsep=2pt,
    parsep=0pt,
    topsep=4pt,
    partopsep=0pt
}
\setlist[enumerate]{
    leftmargin=2.8em,
    itemsep=2pt,
    parsep=0pt,
    topsep=4pt,
    partopsep=0pt
}
\setlist[description]{
    leftmargin=2.8em,
    itemsep=2pt,
    parsep=0pt,
    topsep=4pt,
    partopsep=0pt
}

% 5. กล่องเน้นข้อความ (Callout Box) ไม่ให้ย่อหน้าแรกเยื้องผิดตำแหน่ง
\tcbset{
    enhanced,
    boxrule=0.7pt,
    colback=gray!5,
    colframe=gray!50,
    arc=3pt,
    left=12pt,
    right=12pt,
    top=8pt,
    bottom=8pt,
    before skip=10pt,
    after skip=10pt
}
